\section{Conclusions}
\label{sec:conclusions}

\begin{enumerate}[leftmargin=1.6em,itemsep=4pt,topsep=4pt,parsep=0pt]

\item \emph{What is new.} In the millimetre band the record for nearby stars
is two measurements of one spectral line: \citet{Steffes1994} and
\citet{Mauersberger1996} each covered no more than \SbPriorBandMHz\,MHz at
narrowband resolution around the \SbPriorFreqGHz\,GHz positronium line, with
no search over drift rate and
no epoch-to-epoch test of persistence. This survey covers at
least \SbPriorRatio{} times that bandwidth, searches a grid of drift rates,
and asks of every event whether it returns (\S\ref{sec:background}).

\item \emph{What was searched.} The carrier search is \NWinA{} fine-channel
windows toward \NSysClassA{} systems within 40\,pc, covering \CvUnionA\,GHz
of unique frequency in \CovNIslands{} disjoint intervals. In \EdrNWinA{} of
those windows the track is long enough for the searched drift to cross a
channel, covering \EdrUnionGHz\,GHz toward \EdrNSys{} systems. A coarse-channel
search for a steady unresolved excess adds \NWinB{} windows toward
\PrSysClassB{} systems. Both come from one pass through \NEB{} public ALMA
execution blocks and together span \SurvFreqLo--\SurvFreqHi\,GHz
(\S\ref{sec:sample}).

\item \emph{How deep.} The median system's best window would have produced a
threshold crossing from a carrier of $\EirpNinetySysMedA$\,W nine times in
ten, and \NSysEfifteen{} systems reach $10^{15}$\,W on that measure. Three
things qualify it. It is the completeness of one observation and not of a
confirmation, which needs a second and separated epoch
(\S\ref{sec:confirmability}); it is an estimate rather than a measurement
for \SmcNEstSys{} of the \NSysClassA{} systems, whose best window was not
itself injected into; and it leaves out atmospheric decorrelation, which a
real carrier suffers and an injected one does not, against which a
conservative term of comparable size runs the other way, so the net bias is
of either sign by window (\S\ref{sec:pninetysel}).

\item \emph{What was found.} The search produced \EvNCrossRaw{} threshold
crossings: \CkNCrossCenA{} Class~A in the primary census, \PrFlagXrossWord{}
Class~B, and \CeNTailCrossWord{} Class~A in archival-tail blocks. The
\PrFlagXrossWord{} Class~B crossings are the whole of that search's crossing
population and all come from the single execution block that fails the
data-quality criterion of Appendix~\ref{sec:blockquality}, so the
quality-passing data hold \EvNCrossQp{} crossings, all of them from the
carrier search. \EvNAttrQp{} of those are attributed to
a catalogued molecular transition at the star's own velocity and
\EvNUnattrQp{} are not.
An attribution here is a coincidence of frequency and velocity and not an
identification: \LnAttrSupp{} of the \EvNAttrQp{} carry evidence beyond the
frequency match -- a second transition of the same species, catalogued gas
on the sightline, or the crossing's own recurrence -- and
\LnAttrCoincWord{} rest on the coincidence alone. The unattributed count is
what these fields' own control positions predict, \CeExp{} expected against
the \CeObs{} found, with probability \CeProbNull. That is consistency and
not a demonstration that the crossings are noise: a further \CeAddLim{}
would still have been consistent
(\S\ref{sec:technosearch}, \S\ref{sec:trigunif}).

\item \emph{What the repeat observations exclude.} A transmitter on an
orbiting body arrives at a frequency its own motion has moved, so every
unattributed crossing another block covers was searched again across that
whole window at every drift of the grid. \WdNExc{} of the \WdNCov{} covering
windows hold a statistic above the trigger, which is the number chance
predicts in trials of that size, \WdExpect. Each of those is itself a
crossing, so each of the \TjNPair{} pairs one carrier might have produced
was asked whether a single trajectory fits both. The \TjNPower{} whose
fitted drifts the grid resolves all fail a constant-acceleration track, but
the exclusion claimed is the weaker and safer one -- that no trajectory
inside the searched grid fits both -- and that separates only \TjNExcl.
The remaining \TjNOpen{} cannot be separated -- \TjNPairEdge{} of them
marked \emph{edge}, where one fitted rate is only a lower bound -- and are
\emph{open}: neither a confirmation nor an exclusion, with only the
agreement of their number with chance bearing on them. Read instead at the single cell
its own discovery predicts -- a narrower question, which asks a carrier to
return to that cell and radiate through both observations, and not a test of a
persistent transmitter of any kind -- no crossing returned: the largest
statistic at any predicted cell is \RcTMax{} against a trigger of
\StageNullTrig. The same reading recovers \MkAttrRecur{} of the
\MkAttrTested{} line-attributed crossings it was tested on, each in a
different observing session more than a day away
(Table~\ref{tab:confirm}); what those recoveries establish divides with the
width of the line they fall on, so the registration of a single cell is
tested by the \RgRegNUnres{} readings of \RgRegStar{} alone, and the rest
test the frame correction only to the width of a resolved line. Each
exclusion's depth was measured by injection as well as calculated, the
measured one is the weaker and is the one quoted, and at \HjEtaStar{} the
injection is recovered \HjEtaMeas{} of the time where the calculation gives
one, so that crossing is reported as untested
(\S\ref{sec:etacrv}).

\item \emph{What could have been confirmed at all.} A search limit and a
confirmation limit are different statements. Of the \NSysClassA{} systems the
carrier search covered, \ZcConfSysDay{} were observed again in another
observing session more than a day later at a frequency the first block
covered -- counting every
block the recurrence test reads, not the census alone -- and \ZcConfSysYr{}
more than a year later. For the remaining \ZcNoRepSys{} no repeat exists, so
no transmitted power would have made a confirmation possible
(\S\ref{sec:confirmability}).

\item \emph{Stacking the epochs.} Combining each star's observations in its
own rest frame deepens the flux limit by a median factor \SyStkGainMed. The
deepest result is Proxima Centauri's: $\SzProxEirpCensus$\,W from the census
blocks and $\StkProxEirpCal$\,W once the rest of the archive is added, in
coarse \ChanBModeMHz\,MHz channels, so it bounds a steady unresolved excess
and not a drifting carrier (\S\ref{sec:stack}). Stacking buys that depth
with a tolerance of one channel, and \PxPlanet's own orbit carries a
transmitter across \PxNChanWalk{} of them, so the stacked figure applies to
the one body in that system only if the orbit is within
\PxIncMaxDeg$^\circ$ of face-on. What applies at any inclination is the
single-epoch limit, $\PxEirpSingle$\,W.

\item \emph{What this does not constrain.} These limits hold only inside the
domain actually searched. The drift grid stops at \DriftCeilLo\,nHz,
\AccelSheikhPct{} per cent of the ceiling recommended for a general carrier
search \citep{Sheikh2019}, so an unbound transmitter, one on a rapidly
rotating planet, or one whose orbital and rotational accelerations together
exceed that is not constrained here at all (\S\ref{sec:excludes}). What
transfers beyond this archive is the obstacle: channelisation chosen by other
observers, and not collecting area, decides which search is possible.

\end{enumerate}
